\subsection{Type 1: Independent warping}

The most general approach to modeling warping is to take the $n_w$ warping amplitudes $\bm{\alpha}$ as independent configuration variables.
This generalizes the approach of \cite{reissner1952nonuniform} for torsional warping.

\subsubsection{Equilibrium}

\vspace{1ex}
\begin{ambox}[Type 1: Independent warping]{box:type1-warping}
A Type~1 constitutive response maps the $6+2n_w$ strains $(\bm{\gamma},\bm{\kappa},\bm{\alpha}',\bm{\alpha})$ to the $6+2n_w$ generalized stresses $(\bm{n},\bm{m},\bm{w},\bm{v})$,
\[
  [\bm{n},\bm{m},\bm{w},\bm{v}]
    = \mathcal{S}(\bm{\gamma},\bm{\kappa},\bm{\alpha}',\bm{\alpha}).
\]
The (linearized) momentum balance can be written as
\begin{equation*}
  \begin{aligned}
    \mathbb{B}_{\mathrm{p}}\,\bm{A}_*
      \begin{pmatrix} \bm{n} \\ \bm{m} \end{pmatrix}
      + \bm{f}_{\mathrm{p}}
      &= \dot{\bm{\mu}},
    \\
    \bm{w}' - \bm{v} + \bm{f}_{\mathrm{w}}
      &= \mathbf{0},
  \end{aligned}
\end{equation*}
where $\bm{A}_*$ accounts for the transformation from material to spatial resultants, and $\bm{f}_{\mathrm{p}}$ collects the primary body loads.
\end{ambox}

\subsection{Type 2: Constrained warping}

It is common to assume that the amplitudes of certain warping modes are proportional to a component of one of the primary strain variables $\bm{e}_{\mathrm{p}} = (\bm{\gamma},\bm{\kappa})$.
For example, torsional warping is often assumed to be proportional to the rate of twist
$\varkappa \triangleq \bm{\kappa}\cdot\FrameBasis$.
To formalize this, we partition the warping modes into three disjoint index sets
$\mathcal{I}_v$, $\mathcal{I}_w$, and $\mathcal{I}_f$, corresponding to algebraically constrained, differentially constrained, and free modes, respectively.

We introduce projection operators $\mathbf{L}_{\rm vv},\mathbf{L}_{\rm ww}\in\mathbb{R}^{n_w\times n_w}$ onto the subspaces of algebraically and differentially constrained modes, satisfying
\[
  \mathbf{L}_{\rm vv}^2 = \mathbf{L}_{\rm vv}, \qquad
  \mathbf{L}_{\rm ww}^2 = \mathbf{L}_{\rm ww}, \qquad
  \mathbf{L}_{\rm vv}\mathbf{L}_{\rm ww}
    = \mathbf{L}_{\rm ww}\mathbf{L}_{\rm vv}
    = \mathbf{0}.
\]
The constraints are expressed through linear maps
$\mathbf{L}_{\rm vp},\mathbf{L}_{\rm wp}\in\mathbb{R}^{n_w\times 6}$ via
%
\begin{equation}
  \left\{
  \begin{aligned}
    \mathbf{L}_{\rm vv}\,\bm{\alpha} &= \mathbf{L}_{\rm vp}\,\bm{e}_{\mathrm{p}}, \\
    \mathbf{L}_{\rm ww}\,\bm{\alpha}' &= \mathbf{L}_{\rm wp}\,\bm{e}_{\mathrm{p}},
  \end{aligned}
  \right.
  % \label{eq:linear-warp-constraints}
\end{equation}
%
with compatibility conditions
$\mathbf{L}_{\rm vv}\mathbf{L}_{\rm vp} = \mathbf{L}_{\rm vp}$ and
$\mathbf{L}_{\rm ww}\mathbf{L}_{\rm wp} = \mathbf{L}_{\rm wp}$.

The warping amplitude vector is decomposed as
\begin{equation}
  \bm{\alpha}
    = \bar{\bm{\alpha}}
      + \mathbf{L}_{\rm vv}\bm{\alpha}
      + \mathbf{L}_{\rm ww}\bm{\alpha},
  \label{eq:decompose-alpha}
\end{equation}
where the reduced warping vector
$\bar{\bm{\alpha}} \triangleq (\mathbf{I}-\mathbf{L}_{\rm vv}-\mathbf{L}_{\rm ww})\bm{\alpha}$
contains only the free components.

\subsubsection{Algebraic constraints}

\paragraph{Kinematics}
Using the constraints \eqref{eq:linear-warp-constraints}, the full strain vector $\bm{e}$ can be expressed in terms of the primary strains and the reduced warping variables as
\begin{equation}
  \bm{e}
  =
  \begin{pmatrix}
    \bm{\gamma} \\ \bm{\kappa} \\ \bm{\alpha}' \\ \bm{\alpha}
  \end{pmatrix}
  =
  \underbrace{
  \begin{bmatrix}
    \mathbf{1} & \mathbf{0} & \mathbf{0} & \mathbf{0} \\
    \mathbf{0} & \mathbf{1} & \mathbf{0} & \mathbf{0} \\
    \mathbf{L}_{\rm wn} & \mathbf{L}_{\rm wm} & \mathbf{1}-\mathbf{L}_{\rm ww} & \mathbf{0} \\
    \mathbf{L}_{\rm vn} & \mathbf{L}_{\rm vm} & \mathbf{0} & \mathbf{1}-\mathbf{L}_{\rm vv}
  \end{bmatrix}
  }_{\mathbf{L}}
  \begin{pmatrix}
    \bm{\gamma} \\ \bm{\kappa} \\ \bar{\bm{\alpha}}' \\ \bar{\bm{\alpha}}
  \end{pmatrix},
  \label{eq:warping-projection}
\end{equation}
where $\mathbf{L}_{\rm vn}$ and $\mathbf{L}_{\rm vm}$ are the partitions of $\mathbf{L}_{\rm vp}$ corresponding to $\bm{\gamma}$ and $\bm{\kappa}$, respectively, and similarly for $\mathbf{L}_{\rm wn}$ and $\mathbf{L}_{\rm wm}$.

\begin{remark}
Here $\bar{\bm{\alpha}}$ collects the unconstrained warping amplitudes, while $\bm{\alpha}$ is the union of these free amplitudes with the constrained ones.
For example, given a two-mode warping model $\bm{\alpha} = (\alpha_1,\alpha_2)$ and the constraint
$\alpha_1 = \mathbf{1}\cdot\bm{\kappa}$, we may take
\[
  \mathbf{L}_{\rm vv}
  =
  \begin{bmatrix}
    1 & 0 \\[0.3ex]
    0 & 0
  \end{bmatrix},
  \qquad
  \bar{\bm{\alpha}}
    = \bm{\alpha} - \mathbf{L}_{\rm vv}\bm{\alpha}
    = (0,\alpha_2),
  \qquad
  \bm{\alpha}
    = \bar{\bm{\alpha}} + (\alpha_1,0).
\]
\end{remark}

\paragraph{Equilibrium}

To derive the equilibrium equations for the constrained system we consider the Galerkin projection of the equilibrium residual.
The weak contribution from the warping part is
\begin{equation*}
  \mathcal{G}_{\chi}[\bm{\eta}]
    = \bigl\langle \mathbb{B}_{\mathrm{p}}
        \begin{pmatrix} \bm{n} \\ \bm{m} \end{pmatrix},
        \bm{\eta}_{\mathrm{p}}\bigr\rangle
      - \bigl\langle \bm{w}' - \bm{v} ,\, \delta\bm{\alpha} \bigr\rangle .
\end{equation*}
Using the decomposition \eqref{eq:decompose-alpha} and the constraint equations
\eqref{eq:linear-warp-constraints}, the test functions $\delta\bm{\alpha}$ can be expressed in terms of variations of the reduced variables.
Exploiting the linearity of the inner product and the transpose property of the constraint operators, one obtains
\begin{equation*}
  \mathcal{G}_{\chi}[\bm{\eta}]
    =
    \bigl\langle
      \begin{pmatrix} \bm{n} \\ \bm{m} \end{pmatrix}
      + \mathbf{L}_{\rm vp}^{\top}\bm{v}
      - \mathbf{L}_{\rm vp}^{\top}\bm{w}' ,
      \delta\bm{e}_{\mathrm{p}}
    \bigr\rangle
    - \bigl\langle \bm{w}' - \bm{v} ,\, \delta\bar{\bm{\alpha}} \bigr\rangle .
\end{equation*}

This motivates the definition of augmented stress resultants
\begin{equation}
  \begin{aligned}
    \bar{\bm{n}} &\triangleq \bm{n} + \mathbf{L}_{\rm vn}^{\top}\bm{v}, \\
    \bar{\bm{m}} &\triangleq \bm{m} + \mathbf{L}_{\rm vm}^{\top}\bm{v}, \\
    \bar{\bm{w}} &\triangleq (\mathbf{1}-\mathbf{L}_{\rm vv}-\mathbf{L}_{\rm ww})^{\top}\bm{w}, \\
    \bar{\bm{v}} &\triangleq (\mathbf{1}-\mathbf{L}_{\rm vv}-\mathbf{L}_{\rm ww})^{\top}\bm{v},
  \end{aligned}
  % \label{eq:augment-stress}
\end{equation}
and the corresponding resultant transformation
\[
  \bar{\bm{s}}
  =
  \begin{pmatrix}
    \bar{\bm{n}} \\ \bar{\bm{m}} \\ \bar{\bm{w}} \\ \bar{\bm{v}}
  \end{pmatrix}
  =
  \underbrace{
  \begin{bmatrix}
    \mathbf{1} & \mathbf{0} & \mathbf{0} & \mathbf{L}_{\rm vn}^{\top} \\
    \mathbf{0} & \mathbf{1} & \mathbf{0} & \mathbf{L}_{\rm vm}^{\top} \\
    \mathbf{0} & \mathbf{0} & \mathbf{1}-\mathbf{L}_{\rm vv}^{\top}-\mathbf{L}_{\rm ww}^{\top} & \mathbf{0} \\
    \mathbf{0} & \mathbf{0} & \mathbf{0} & \mathbf{1}-\mathbf{L}_{\rm vv}^{\top}-\mathbf{L}_{\rm ww}^{\top}
  \end{bmatrix}
  }_{\mathbf{L}^*}
  \begin{pmatrix}
    \bm{n} \\ \bm{m} \\ \bm{w} \\ \bm{v}
  \end{pmatrix}.
\]

\vspace{1ex}
\begin{ambox}[Type 2: Constrained warping]{box:warp2}
For linear constraints on warping of the form \eqref{eq:linear-warp-constraints}, a Type~1 constitutive response can be re-used by the transformation
\[
  \bar{\bm{s}} = \mathbf{L}^*\,\mathcal{S}(\mathbf{L}\,\bar{\bm{e}}),
\]
where $\bar{\bm{s}} = (\bar{\bm{n}},\bar{\bm{m}},\bar{\bm{w}},\bar{\bm{v}})$ and
$\bar{\bm{e}} = (\bm{\gamma},\bm{\kappa},\bar{\bm{\alpha}}',\bar{\bm{\alpha}})$.
The momentum balance becomes
\begin{equation}
  \begin{aligned}
    \mathbb{B}_{\mathrm{p}}\,\bm{A}_*
      \bigl[\bar{\bm{s}}_{\mathrm{p}}
            - \mathbf{L}_{\rm vp}^{\top}\bm{w}'\bigr]
      + \bm{f}_{\mathrm{p}}
      &= \dot{\bm{\mu}},
    \\
    \bar{\bm{w}}' - \bar{\bm{v}} + \bm{f}_{\mathrm{w}}
      &= \mathbf{0},
  \end{aligned}
  \label{eq:type2-equilibrium}
\end{equation}
where $\bar{\bm{s}}_{\mathrm{p}} = (\bar{\bm{n}},\bar{\bm{m}})$ denotes the primary resultants.
\end{ambox}

\begin{remark}
The treatment of the $\mathbf{L}_{\rm vp}^{\top}\bm{w}'$ term in \eqref{eq:type2-equilibrium} reveals a fundamental distinction between section-level and element-level operations.
The constitutive evaluation at a section is pointwise in $\reflenx$ and therefore cannot access spatial derivatives.
Consequently:
\begin{itemize}
  \item The \emph{section wrapper} receives $(\bm{\gamma},\bm{\kappa},\bar{\bm{\alpha}}',\bar{\bm{\alpha}})$, expands them via $\mathbf{L}$ to obtain the full strain $\bm{e}$, evaluates the Type~1 constitutive response $\mathcal{S}(\bm{e})$, and returns the partially augmented resultants
  \[
    \bar{\bm{n}} = \bm{n} + \mathbf{L}_{\rm vn}^{\top}\bm{v}, \qquad
    \bar{\bm{m}} = \bm{m} + \mathbf{L}_{\rm vm}^{\top}\bm{v},
  \]
  together with the full $\bm{w}$-field.
  \item The finite element must then compute $\bm{w}'$ from its interpolated bimoment field and apply the correction
  $-\mathbf{L}_{\rm vp}^{\top}\bm{w}'$ to the primary equilibrium equations.
\end{itemize}
This division of labor implies that Type~2 elements require higher-order interpolation to represent $\bm{w}'$ accurately, making them more complex to implement than standard beam elements.
This motivates the Type~3 formulation, in which the assumption $\bar{\bm{\alpha}}'=\mathbf{0}$ (uniform warping) eliminates the $\bm{w}'$ term entirely, keeps all manipulations at the section level, and allows standard $6$-DOF beam elements to be used without modification.
\end{remark}
